% ED26-SRC-036-010 continuation, Struss 2008 sections 10.6--10.7.
\chapter{Model-Based Remedies, Configuration and Failure Analysis}
\label{sec:p16-model-based-remedy-and-configuration}

A diagnosis describes what may be wrong; it is not yet an action choice.
Struss's remedy development asks what should be restored, what still persists
after intervention, which remaining diagnoses must be accommodated, and
whether a proposed target can be reached safely
\cite[secs.~10.6--10.7, pp.~446--458]{Struss2008ModelBasedProblemSolving}.
The component/process/model construction is in
\cref{sec:p16-qualitative-compositional-model-formulation}.
Here we construct the additional target and execution reasoning, rather
than renaming the static diagnosis or a statistical intervention.

The chapter's simplified action-history definition assumes instantaneous
successful actions. Its cited general repair formalization allows action
failure, but those are distinct versions. The finite constructions below
state their transition, observation and coverage conditions; they are not
claims that the chapter supplied a complete planner for every physical system.

\section{From a diagnostic hypothesis to a reconfiguration target}
\label{sec:p16-rem-targets}

Replacement usually changes a component mode; switching-based reconfiguration
changes operating states while faults remain. Changing a set point changes
an operating region. Switching control software can avoid a suspect sensor.
Introducing a cable, heater or oxidizer changes structural/model content.
Even an intrusive test can change the system. Struss's action and goal
histories therefore augment, rather than merely re-run, diagnosis histories.
Intermediate goals include avoiding line overload and preserving power to
critical consumers \cite[pp.~447--450]{Struss2008ModelBasedProblemSolving}.

\begin{hbdefinition}{def:p16-rem-target-input}{Native state assignments and persistent data}
Let \(SD=LIB\cup STRUCTURE\) be the source's system description,
and \(M_1,\ldots,M_k\), \(k\ge1\), remaining complete diagnostic
mode assignments. A finite set \(C_R\) of reconfigurable components has
nonempty finite state domains \(S_c\). A complete target assignment
\(s\in\prod_{c\in C_R}S_c\) also denotes the conjunction of its
chosen state atoms. Keep fault modes fixed for this target-selection task.
For each hypothesis let \(O_{P,i}\) be the correctly determined observations
which persist after the proposed type of intervention. Put
\(W_i=SD\cup\{M_i\}\cup O_{P,i}\); goals \(G_i\) are finite
conjunctions, or finite theories identified with their conjunctions.

A common consistency target has \(W_i\cup\{s\}\cup G_i\) consistent
for every \(i\). A common guaranteed target has \(W_i\cup\{s\}\)
consistent and entailing \(G_i\) for every \(i\). ``For every model''
does not mean that each model may choose a different target state.
\end{hbdefinition}

Opening a breaker changes consumer voltages, but not necessarily generator
availability. Carrying old consumer voltage into the target can falsely
reject restoration; discarding the inactive-generator observation can falsely
accept it. Persistent data is a supplied/established frame condition, not
whatever facts make the target consistent. A weak open-line model which
leaves voltage free also permits a desired voltage without producing it.
Source-connected supply must be modeled, not inferred from the goal sentence
\cite[pp.~450--452]{Struss2008ModelBasedProblemSolving}.

\begin{hbdefinition}{def:p16-rem-change}{Actual changes, common entailed states and cost}
With known actual assignment \(s_A\), the change set is
\(\Delta(s)=\{c:s(c)\ne s_A(c)\}\).
With incomplete information, let \(s_{M,i}\) contain all component-state
atoms entailed by \(SD\cup\{M_i\}\cup O\), and let
\(s_{common}=\bigcap_i s_{M,i}\). Struss's generalized change measure
is the set of target atoms \(s\setminus s_{common}\).
It counts a specified target whose old state is unknown, even if one
possible actual state already equals it. It is not the same as the actual
number of physical changes. Inclusion minimality is different from weighted
cost, irreversible-actuation risk or downtime.
\end{hbdefinition}

For example, if a valve's actual state may be OPEN or CLOSED but no state
is entailed, an OPEN target appears in the generalized change set. In the
OPEN possible world it needs no state change, while in the CLOSED world
it does. A command which sets OPEN idempotently may handle both; an action
whose precondition is CLOSED does not. Broken components are not controllable
merely because their state names appear in \(C_R\).

\begin{hbmethod}{def:p16-rem-target-search}{Context-tagged state-conflict search}
Generate complete targets without repetition, using the stated cost/order
criterion. Before composing a target, reject it if it extends a cached
inconsistent partial state assignment from one of the current contexts
\(W_i,G_i\). Otherwise check each composed target model for consistency.
For guaranteed targets, also check whether a model satisfying the negation
of the conjunctive goal exists. Cache supporting state conflicts when
available. At least one assigned state in such a conflict must change
in an acceptable target.
Return feasible common targets, then remove those dominated under the chosen
change criterion. Conflicts from old persistent data or a revised structure
cannot be reused without re-establishing their meaning.
Failure to entail a goal alone is not an inconsistent partial assignment:
changing another state may remove its countermodel.
\end{hbmethod}

\begin{hbtheorem}{thm:p16-rem-target-search}{proved=annex:p16.rem-target-search}{Full finite target coverage and its oracle premise}
Let \(N=\prod_c|S_c|\) and supply total exact consistency/consequence
checks. The target construction terminates after at most \(kN\)
consistency checks for common possibility, or \(2kN\) for the stated
guarantee tests, and returns exactly its common targets before minimization.
Sound context-matched partial-state conflicts remove no acceptable target.
The minimal-change filter returns exactly its inclusion-minimal change sets.
With finite nondecreasing cost on such sets, a cheapest feasible target has
a minimal-change feasible target of no greater cost. This does not prove
a plan exists, nor prove termination of unrestricted first-order inference.
\end{hbtheorem}

The source compares state-conflict search to fault identification, not to
constraint suspension with unconstrained faults. Closing one switch can
disconnect another load or cause an overload; every revised target must
be recomposed. There is no universal OK-versus-faulty preference among
operating states. Generator start-up and switch operations can have very
different costs \cite[pp.~452--453]{Struss2008ModelBasedProblemSolving}.

\section{Action histories require guards, effects and frames}
\label{sec:p16-rem-execution}

Struss's source criterion joins \(SD\), a mode history \(MH\), an
action history \(AH\), observations \(OBSH\) and goals \(GOALSH\).
Consistency permits a goal-satisfying history; explanatory achievement
requires the action/history model to entail the goals. Applicability and
consistency must also hold, so an impossible action cannot achieve every
goal by vacuous entailment. The goal history may constrain intermediate
intervals as well as the final one
\cite[Definition~10.18, p.~449]{Struss2008ModelBasedProblemSolving}.

\begin{hbdefinition}{def:p16-rem-execution-model}{A bounded component-execution realization}
For each remaining physical/model hypothesis \(i\), construct a finite
state set \(X_i\) from admissible component modes, operating states,
persistent parameters and any bounded structural choices. Include all
memory relevant to future behavior. A repair may change a current mode;
the hypothesis tag does not freeze that mode after replacement.
Let \(\mathcal X=\coprod_i\{i\}\times X_i\), with nonempty initial
possibility set \(B_0\subseteq\mathcal X\) selected by the actual data.
The distinct tags are retained even when two states have equal observables.

For each allowed action \(a\), its relation \(T_a\subseteq\mathcal X^2\)
must be generated from the library, preconditions, declared effects,
frame conditions and the post-action structural/behavior constraints.
It preserves the model tag. On every state where its guard holds it has
at least one successor; otherwise it has none. Safety and terminal goals
are subsets \(\mathcal S,\mathcal G\subseteq\mathcal X\).
Between-action requirements must be encoded or proved separately: this
finite realization does not silently establish continuous physical safety.
\end{hbdefinition}

For a healthy controllable switch, an idempotent SetSwitch action sets its
operating state, preserves its mode and other unaffected switches, and
recomputes voltages from the post-state network. Derived voltage is not
framed as unchanged. Replacing a component can set its mode to normal but
must use the appropriate replacement type/parameters; if structure changes,
the post-state model is recomposed. Reset success or continued control of
a faulty actuator is not inferred from the action's name.

\begin{hbdefinition}{def:p16-rem-strong-plan}{All modeled executions, not one optimistic witness}
For \(B\subseteq\mathcal X\), define
\(Post(B,a)=\{y:\exists x\in B,(x,y)\in T_a\}\).
Action \(a\) is strongly applicable when every \(x\in B\) has a
successor. A strong plan uses strongly applicable actions, has
\(B_0\subseteq\mathcal S\) and every successive \(Post\) set inside
\(\mathcal S\), and ends with a nonempty set inside \(\mathcal G\).
A weak history needs only one admissible goal-reaching execution.
An action cannot silently discard inconvenient initial worlds by having
no successor for them.
\end{hbdefinition}

\begin{hbmethod}{def:p16-rem-plan-search}{Complete finite strong open-loop planning}
If \(B_0\not\subseteq\mathcal S\), there is no plan satisfying the
stated safety requirement from its initial state. Otherwise start from
\(B_0\). For each reached nonempty safe possibility set, generate
each strongly applicable action whose complete successor set is safe.
Store the successor set and an incoming action path, reusing a previously
stored set. Breadth-first expansion returns the first set contained in
the terminal goal, or exhausts the finite graph. When future rules or
goal phases depend on past events, include that information in the state
before sharing nodes; belief equality alone then means a sufficient state.
\end{hbmethod}

\begin{hbtheorem}{thm:p16-rem-plan-search}{proved=annex:p16.rem-plan-search}{Prefix coverage, soundness and finite strong-plan completeness}
For the declared finite execution model, iterated \(Post\) gives exactly
all states reachable after each action prefix from any initial hypothesis,
with all modeled alternatives retained. If \(|\mathcal X|=n\) and
there are \(m\) actions, the search has at most \(2^n-1\) nonempty
nodes and \(m(2^n-1)\) candidate action expansions. It terminates,
returns a strong plan exactly when one exists, and breadth-first search
returns one with minimum action count. This is not minimum monetary cost
or a completeness theorem for an unbounded structure-changing domain.
If the initial actual world is represented, model applicability guarantees
physical applicability, all actual action outcomes are represented by the
relations, and decoded safety/goals imply physical safety/goals, the returned
strong plan also gives that physical guarantee. These coverage and guard
conditions are additional, not consequences of graph search.
\end{hbtheorem}

\section{The power-network cases: restoration, ordering and incomplete diagnosis}
\label{sec:p16-rem-power}

\begin{hbdefinition}{def:p16-rem-power-transfer}{A bounded realization of the source's intermediate-goal case}
Two Boolean route switches \(x,y\) connect a critical consumer; exactly
one may be closed without overload. They are healthy controllable switches
fed by a known active source. In this idealized switching model,
critical supply is \(x\lor y\). Each action sets exactly one switch,
is instantaneous and successful, and frames the other. The initial state
is \((1,0)\); the required target is \((0,1)\). Safety requires
\(x\lor y\) and \(\neg(x\land y)\) at every execution state.

For a second case add an isolated, adequately rated temporary backup switch
\(b\). Supply becomes \(x\lor y\lor b\), while the main-route
overload condition remains \(\neg(x\land y)\). Backup actuation requires
known healthy available backup power. Initially and finally \(b=0\).
These isolation/availability assumptions are part of the model, not consequences
of calling a device a backup.
\end{hbdefinition}

\begin{hbtheorem}{thm:p16-rem-power-transfer}{proved=annex:p16.rem-power-transfer}{A feasible target with no safe plan, and the required temporary actions}
Both two-switch endpoints are feasible and safe, but no
sequence of the stated single-switch actions connects them safely.
With the stated backup available, the sequence
\[
 b:=1,\quad x:=0,\quad y:=1,\quad b:=0
\]
is safe and reaches the target. Four actions are necessary in this model.
The target changes only the two main states; the temporary backup actions
are not visible in that target difference. Simultaneous double-switch
actuation or relaxed continuity of supply would change the first conclusion.
\end{hbtheorem}

This spells out Struss's distinction between computing a goal state and
planning its realization. Its assumptions are an explicit small instance,
not an undocumented numerical model of an actual power grid
\cite[p.~453]{Struss2008ModelBasedProblemSolving}.
If backup availability is unknown, the EnableBackup guard fails in the
unavailable possible world. It cannot be treated as a strong action by
discarding that world. A measurement may reveal availability, but an
unavailable branch still needs a different remedy, a relaxed goal, or a
structural extension; it is not cured by diagnosis alone.

The source also permits restoration before exact fault localization.
For three isolated routes, suppose two diagnostic hypotheses respectively
identify route A or route C as unavailable, while route B and its controllable
isolators are healthy in both. With a known active generator, disconnecting
A and C and connecting B supplies the consumer under either hypothesis.
This conclusion follows by recomposing both models: neither target relies
on a suspect route. It fails if the generator is inactive, B is also suspect,
or an unmodeled bridge defeats the assumed isolation. The common remedy
is not evidence that one of the initial diagnoses was uniquely true
\cite[p.~452]{Struss2008ModelBasedProblemSolving}.

\section{Tests and remedies in one observation-guided process}
\label{sec:p16-rem-policy}

When the next action depends on a test result, a single open-loop sequence
is insufficient. Choose a finite observation map \(h:\mathcal X\to Z\).
After action \(a\), its observable result \(z\) gives the nonempty
branch \(B_{a,z}=Post(B,a)\cap h^{-1}(z)\).
Only observed results license branching, not the hidden model tag.
An intrusive test must use its actual transition/model revision before
this filtering, rather than pretending it observes the unchanged device.

\begin{hbmethod}{def:p16-rem-policy-search}{A finite strong observation-policy construction}
Let \(Win_0\) be the nonempty safe belief sets contained in the goal.
Repeatedly add a nonempty safe belief set \(B\) if an action is strongly applicable,
has a safe complete successor set, and every nonempty observation branch
belongs to the previous winning set. Store that action and its branches.
Stop when no set is added. The insertion round gives a decreasing execution
rank; at \(Win_0\), stop acting. Goal-phase or history information belongs
in the state before this construction is used.
\end{hbmethod}

\begin{hbtheorem}{thm:p16-rem-policy-search}{proved=annex:p16.rem-policy-search}{Exactly the finite-depth strong policies at this scope}
The construction terminates after at most \(2^{|\mathcal X|}-1\)
strict addition rounds. Membership in \(Win_j\) is equivalent to having
an observation-guided strong policy of depth at most \(j\). The final
set therefore contains exactly the beliefs with finite-depth strong policies.
All possible observation branches must succeed. Almost-sure termination,
fairness-dependent cycles and arbitrary unbounded histories are not identified
with this criterion.
\end{hbtheorem}

If an action may fail without changing state, that branch has to remain
represented. Repeating the action can be plausible under a probabilistic
or fairness premise, but a sequence of endlessly failing outcomes is still
possible without it. The source's simplified instantaneous-successful-action
definition and its referral to failure-aware repair plans must not be merged
without those new premises \cite[pp.~447--449]{Struss2008ModelBasedProblemSolving}.

\section{Located process therapy preserves amounts, not old derivatives}
\label{sec:p16-rem-therapy}

Struss's iron-water case adds an oxidizer object at a stated location,
activating oxidation and producing iron oxide. The therapy introducibles
are action triggers, not unexplained disturbance objects. Its consistency
criterion joins DEC, the persistent situation, LIB and therapy goals.
It does not supply a multistep intervention planner by itself
\cite[Definition~10.21, pp.~453--454]{Struss2008ModelBasedProblemSolving}.

\begin{hbdefinition}{def:p16-rem-therapy-model}{A complete-influence tank instance with separate before/after times}
In this bounded constant-volume tank model, dissolved-iron concentration
is \(q\), with positive current deviation above threshold \(H\).
At an instantaneous intervention, \(q_1=q_0>H\). A retained introduction
process contributes \(p\ge0\) to its derivative. A TankOxidizer trigger
adds a correctly located oxidizer and an oxidation process contributing
\(-k\), so \(q'_1=p-k\); these are all derivative influences.
Supply \(Q=Q_0>0\) is retained. Rates independently range over nonempty
closed intervals \([p_{min},p_{max}]\), \([k_{min},k_{max}]\).
Both lower bounds are nonnegative and both upper bounds finite.
The immediate therapy goals are \(q'_1<0\) and continued supply.
The old positive derivative is indexed at time zero, not framed as
\(q'_1=q'_0\). This is an idealized process-library instance, not a
chemical dosing prescription.
\end{hbdefinition}

\begin{hbtheorem}{thm:p16-rem-therapy}{proved=annex:p16.rem-therapy}{The net effect, instantaneous limit and service-goal rejection}
Negative post-intervention derivative is guaranteed exactly when
\(k_{min}>p_{max}\), and possible exactly when \(k_{max}>p_{min}\).
The concentration cannot already satisfy \(q_1\le H\) at that instant.
A StopPump trigger setting \(Q_1=0\) violates continued supply even if
it reduces iron delivery. An unmodeled additional positive influence can
invalidate the derivative guarantee; a ReservoirOxidizer trigger does not
inherit the tank's immediate removal term without its transport model.
\end{hbtheorem}

Retaining amounts but recomputing derivatives is the native persistence
distinction. Keeping an old derivative observation as an unindexed post-state
fact would contradict the proposed therapy artificially. Goals not currently
violated still matter: stopping city supply is not a satisfactory way to
remove an unwanted delivered substance. Minimal DEC is a static inclusion
criterion; timing, duration, side effects, dose constraints and transport
can require more model content and a sequence of interventions.

\section{Configuration is a bounded design search, not fault localization}
\label{sec:p16-rem-configuration}

Struss's configuration object is \(CONF=STRUCTURE\cup PARS\),
with \(DD\cup CONF\cup CONFGOALS\) consistent. DD can leave
structure unspecified: the chosen components, connections or parameters
are the result, not a fixed observed device whose modes must be diagnosed.
Variant design starts from a predecessor with a bounded modification set;
unrestricted innovative design can have an infinite structural search space
\cite[Definition~10.22, pp.~455--456]{Struss2008ModelBasedProblemSolving}.

The source's telephone-switching referral uses dependencies among types,
parameters and viable structures, not necessarily full physical behavior
models. The following finite capacity/power instance makes that distinction
executable without pretending to reconstruct the entire cited industrial
system or inventing its published parameters.

\begin{hbdefinition}{def:p16-rem-card-configuration}{An explicit finite configuration catalog}
Two identical slots each choose Empty, type A or type B. Type A supplies
four capacity units and uses three power units; B supplies six and uses
five. Let \(n_A,n_B\) count the cards. DD imposes the type/count semantics
and \(n_A+n_B\le2\); goals require capacity at least ten and power
at most eight:
\[
 4n_A+6n_B\ge10,\qquad3n_A+5n_B\le8.
\]
There are no additional ordering or compatibility restrictions in this
declared small catalog. These units are stipulated model data, not empirical
specifications for a telephone product.
\end{hbdefinition}

\begin{hbtheorem}{thm:p16-rem-configuration}{proved=annex:p16.rem-configuration}{Every catalog configuration and its actual boundary}
The only feasible count pair is \((n_A,n_B)=(1,1)\), and the two slot
assignments AB and BA are the entire solution set. Removing either card
violates the capacity goal. More generally, finite catalogs and bounded
slot/connection/parameter domains give a finite candidate space; exhaustive
candidate composition with a total exact DD/goal check returns exactly
its configurations, and sound partial design-decision conflicts prune none.
This conclusion does not extend to unbounded numbers of generated components
or prove that consistent catalog constraints guarantee unmodeled functionality.
\end{hbtheorem}

If the predecessor uses two A cards, one type substitution suffices for
the stated configuration goals. Implementing that substitution is a separate
replacement/action task with its own compatibility, data-persistence and
execution conditions. Design conflicts tell which prior choices cannot
remain together; they do not invent a viable initial structural design.
The chapter's VHDL hardware-debugging referral compares specified goals to
a simulated hierarchy, making intended behavior explicit rather than
assuming all implemented components imply it.

\section{FMEA: fault/scenario effects, masking and detectability}
\label{sec:p16-rem-fmea}

For a fixed design, fault assignment \(M_F\) and scenario SCEN, use
\(\Gamma_{F,SCEN}=SD\cup\{M_F\}\cup SCEN\). Require it to
be consistent before classifying effects. A named unwanted effect E is
guaranteed when \(\Gamma\entails E\), and possible when
\(\Gamma\cup\{E\}\) is consistent. The functions/effects and
scenarios are inputs to the analysis. Severity, probability and detectability
are distinct further inputs to criticality; logical possibility supplies
none of their numerical values \cite[p.~456]{Struss2008ModelBasedProblemSolving}.

\begin{hbdefinition}{def:p16-rem-redundancy}{The source's masking issue in a fully specified two-route case}
Two independent isolated supply routes are both commanded on and fed by
a known active source. A normal route supplies \(q_i=1\), an open-fault
route supplies \(q_i=0\), and total service is
\(s=q_1\lor q_2\). Fault sets are subsets of the two routes.
The sole coarse observable is \(s\); the unwanted effect is \(s=0\).
No bridge fault, shared-source failure or cross-route interaction is in
this small library. Adding them would require different models.
\end{hbdefinition}

\begin{hbtheorem}{thm:p16-rem-fmea}{proved=annex:p16.rem-fmea}{All effect cases and the masked fault distinction}
For nonempty models, two exact consistency checks of \(\Gamma\cup\{E\}\)
and \(\Gamma\cup\{\neg E\}\) distinguish guaranteed E, excluded E,
and possible-but-not-guaranteed E. In the two-route case, no faults or
either single open fault exclude loss of service; the double open fault
guarantees it. Neither single fault is detectable from \(s\) alone under
the given scenario because its sole value equals the normal model's value.
Observing \((q_1,q_2)\) distinguishes all four assignments. Thus restricting
FMEA to singles misses the masked double-fault effect.
\end{hbtheorem}

The single/double-fault restriction is a chosen analysis scope, not evidence
that other disturbances are impossible. If the scenario does not activate
a route, that route's open fault can also be unobservable. Conversely an
additional bridge fault can defeat the isolation/OR model even though all
its enumerated cases were analyzed exactly. A library's coverage of relevant
faults is separate from complete enumeration of its selected cases.

\section{Software debugging changes a design, not an old component mode}
\label{sec:p16-rem-debugging}

The source distinguishes a physical device assumed to have worked before
failure from a program whose design may omit a needed computation or contain
wrong structure. Existing code units can execute their specified semantics
while violating the intended function. A well-designed-system assumption
therefore cannot be used to hide the specification. Positive/negative tests
or an abstract specification supply GOALS; language semantics supplies the
connections between program fragments
\cite[pp.~456--458]{Struss2008ModelBasedProblemSolving}.

\begin{hbdefinition}{def:p16-rem-debug-model}{A missing computation and a bounded design revision}
Over integer inputs, the initial program returns \(y=x\). Its intended
specification is \(y=x+1\); the input-zero test expects one but the
actual original output is zero. A finite repair catalog chooses
\(\delta\in\{0,1\}\) and changes the program model to
\(y=x+\delta\). This is a design choice, not an assertion that the
original return instruction has an unforeseen physical failure.
Old measured output belongs to the old version; it is not a persistent
post-revision expected output.
\end{hbdefinition}

\begin{hbtheorem}{thm:p16-rem-debugging}{proved=annex:p16.rem-debugging}{Catalog adequacy, all inputs, and an overfitting counterexample}
The original execution relation and intended output relation are inconsistent
for every shared integer input. In the stated repair catalog, exactly
\(\delta=1\) satisfies the input-zero expectation, and it entails the
intended relation for every integer input. But the program
\(f(0)=1\), \(f(x)=x\) for \(x\ne0\), also passes that one
expectation and fails the specification at input one. Finite-test consistency
is therefore not general design correctness, even when the chosen small
catalog happens to have one adequate surviving revision.
\end{hbtheorem}

An unrestricted faulty-return hypothesis could permit the expected output
without supplying the missing increment, its placement or correct behavior
on other inputs. Code-component blame is not a complete repair construction.
Structural bugs, omitted operations and arbitrary software fault models
remain outside that finite catalog. The source's knowledge-base-debugging
case similarly treats positive and negative examples as goals for a
configuration theory, not observations which prove the theory well-designed.

\begin{hbmethod}{def:p16-rem-consultation}{Completing the source's diagnosis-to-action workflow}
State what functionality must be restored and which diagnoses matter to
the action choice. Identify persistent data, controllable components and
the actual revision type. Construct common target models, then separately
compose guards/effects/frames and plan safe histories. Branch only on
actual observations and retain unknown/failure outcomes. For process therapy,
type the intervention location, new objects/processes, net influences and
post-state goals. For configuration/FMEA/debugging, declare the catalog,
scenario/fault set and intended-function specification rather than relying
on physical-diagnosis normality. Keep model coverage, algorithm completeness,
physical applicability and observed execution success distinct.
\end{hbmethod}

\index{remedy proposal}\index{reconfiguration!planning}\index{configuration}
\index{failure-modes-and-effects analysis}\index{debugging!model based}
